\documentclass[letterpaper,10pt]{article}

% --- Packages ---
\usepackage[empty]{fullpage}
\usepackage{titlesec}
\usepackage[usenames,dvipsnames]{xcolor}
\usepackage{enumitem}
\usepackage[hidelinks]{hyperref}
\usepackage{fancyhdr}
\usepackage{paracol}
\usepackage{tabularx}
\usepackage{fontawesome5}
\usepackage{multicol}
\usepackage{geometry}
\usepackage{parskip}

% --- DOCUMENT SETUP ---
% Set margins for the page
\geometry{
    left=0.5in,
    right=0.5in,
    top=0.6in,
    bottom=0.6in
}

% --- Setup ---
\pagestyle{fancy}
\fancyhf{} % clear all header and footer fields
\renewcommand{\headrulewidth}{0pt}
\renewcommand{\footrulewidth}{0pt}
\fancyfoot[R]{\small\thepage}

\raggedbottom
\raggedright
\setlength{\tabcolsep}{0in}

\titleformat{\section}{\vspace{-5pt}\scshape\raggedright\large}{}{0em}{}[\color{black}\titlerule \vspace{-5pt}]
\titleformat{\subsection}{\vspace{-3pt}\scshape\raggedright\normalsize}{}{0em}{}[\color{black}\vspace{-6pt}]

\begin{document}

% -------------------- HEADING --------------------
\begin{center}
    \textbf{\Huge \scshape Vidhata Jayaraman} \\ \vspace{5pt}
    \small
    % Using a more compact line for contact info
    \faIcon{envelope} \href{mailto:vidhata2@illinois.edu}{vidhata2@illinois.edu} $ \cdot $
    \faIcon{phone-square-alt} 847-420-5370 $ \cdot $
    \faIcon{safari} \href{https://dxdt14.github.io/}{https://dxdt14.github.io/} $ \cdot $
    % \faIcon{github} \href{https://github.com/dxdt14}{\underline{github.com/dxdt14}} $ \cdot $
    \faIcon{linkedin} \href{https://www.linkedin.com/in/vidhata-jayaraman/}{/in/vidhata-jayaraman} $ \cdot $
    \faIcon{book} \href{https://scholar.google.com/citations?user=f16twSMAAAAJ&hl=en}{Google Scholar}
\end{center}

\vspace{-3pt}

\setcolumnwidth{0.5\textwidth, 0.5\textwidth}
\setlength{\columnsep}{1.15cm}
\begin{paracol}{2}
    
    % -------------------- EDUCATION & COURSEWORK --------------------
    \section{Education}
    \textbf{University of Illinois Urbana-Champaign (UIUC)} \\
    B.S. in Computer Engineering $|$ B.S. in Mathematics \\
    Anticipated Graduation: May 2026 | \textit{GPA: 3.97/4.0} \\
    {\small Dean's List (Top 20\% of Student Body)} \\
    
    \vspace{-3pt}
    
    \section{Research Interests}
    Machine Learning \& Optimization, Information Theory \& Quantum Information, High-dimensional Statistics, Geometric Analysis \& Uncertainty Quantification \\
    
    \switchcolumn
    
    \section{Relevant Coursework}
    \textbf{Math}: Real Analysis; Complex Analysis; Abstract Algebra; Linear Algebra; Probability Theory; Algebraic Topology; Graph Theory; Differential Equations; Probability \& Measure \\
    \vspace{10pt}
    \textbf{Applied Math/CS}: Optimization; Machine Learning; Deep Learning for CV; Deep Generative Models; Random Processes; Information Theory; Quantum Information Theory; Algorithms; Data Structures; Signal Processing \\
    
\end{paracol}

\vspace{3pt}

% -------------------- RESEARCH EXPERIENCE --------------------
\section{Research Experience}

\textbf{Research with Professor Lav Varshney at UIUC} \hfill February 2024 -- Present\vspace{-5pt}
\subsection{Current Projects}
\begin{itemize}[itemsep=-5pt, parsep=5pt, leftmargin=0.25in]
    \item \textbf{Senior Thesis}: \textit{Information-theoretic lower bound for Knowledge Distillation in LLMs} (Primary researcher): Attempting to find information-theoretic lower bounds for knowledge distillation with a focus on LLMs
    \item \textit{Extending Community Detection/Graph Clustering analysis to Quantum Networks} (Primary researcher): Extending the "No Free Lunch" Theorem and community detection algorithms from the classical graph setting to quantum networks
\end{itemize}\vspace{-5pt}
\subsection{Past Projects}
\begin{itemize}[itemsep=-5pt, parsep=5pt, leftmargin=0.25in]
    \item \textit{Emergent Capabilities in Transformers}: Experimenting with Modern Hopfield Networks and other Neural Associative Memories to understand emergent capabilities as model size increases and their connection to Transformers
    \item \textit{SwitchCIT: Switching for Continual Instruction Tuning of Large Language Models}: Identified clustering in the final layer of an LLM following continual learning and used this to create a switch network which helps avoid catastrophic forgetting
\end{itemize}

\textbf{Research with Professor Xu Chen at UIUC} \hfill March 2023 -- February 2024\vspace{-5pt}
\begin{itemize}[itemsep=-5pt, parsep=5pt, leftmargin=0.25in]
    \item Created a Physics Informed Neural Network (PINN) to model Voltage and Electric field from a charged circle/sphere inside a grounded box.
    \item Utilized a version of the Deep Galerkin Method (DGM) to estimate the differential equation modeling the system.
    \item Implementation was used \& modified by Samsung engineers for use in their own research and modeling.
    \item Implemented an operator estimator for an RLC (Resistor, Inductor, Capacitor) circuit to model for any R, L, and C.
\end{itemize}

% -------------------- INDUSTRY EXPERIENCE --------------------

\section{Industry Experience}
\textbf{AbbVie} $|$ \textit{Machine Learning Research Intern} \hfill May 2025 -- August 2025\vspace{-5pt}
\begin{itemize}[itemsep=-5pt, parsep=5pt, leftmargin=0.25in]
    \item Implemented the SubgraphX method for GNN explainability, modified for multi-target regression and use with 3D input.
    \item Method correctly identified many important regions in molecules for property prediction, strengthening an internal site-of-metabolism tool.
    \item Conducted novel Uncertainty Quantification research using metric learning and distribution parameter estimation.
\end{itemize}

\textbf{Johns Hopkins University Applied Physics Laboratory} $|$ \textit{Data Science Intern} \hfill June 2024 -- August 2024\vspace{-5pt}
\begin{itemize}[itemsep=-5pt, parsep=5pt, leftmargin=0.25in]
    \item Created a chat-bot with Retrieval Augmented Generation (RAG) for retrieving information in technical documents.
    \item Implemented a Bayesian approach towards reliability growth planning (RGP) using advanced optimization techniques.
    \item Manuscript in preparation to be submitted to IEEE.
\end{itemize}

\textbf{National Institute of Standards and Technology (NIST)} $|$ \textit{Research Intern} \hfill June 2023 -- August 2023\vspace{-5pt}
\begin{itemize}[itemsep=-5pt, parsep=5pt, leftmargin=0.25in]
    \item Created a small GPT-2 model which utilized Cloze Probabilities to identify abnormal sentences in text data.
    \item Demonstrated that analysis of outliers in dimensionally reduced text embeddings can provide similar results, and compared dimensionality reduction methods.
\end{itemize}

\textbf{Brunswick i-JET Research Lab} $|$ \textit{Autonomous Simulation Intern} \hfill January 2023 -- May 2023\vspace{-5pt}
\begin{itemize}[itemsep=-5pt, parsep=5pt, leftmargin=0.25in]
    \item Utilized Robotic Operating System (ROS) to create maps using Simultaneous Localization and Mapping (SLAM).
    \item Utilized theories of fluid dynamics and wake physics to build an autonomous ``perfect" wake generator and used ROS to visualize and analyze data.
\end{itemize}

\pagebreak

\setcolumnwidth{0.45\textwidth, 0.65\textwidth}
\setlength{\columnsep}{1.15cm}
\begin{paracol}{2}
    
    % COLUMN 1
    
    % -------------------- SKILLS --------------------
    \section{Skills \& Expertise}
    \textbf{Deep Learning}: PyTorch, TensorFlow\\
    \textbf{Natural Language}: NLTK, spaCy, LangChain\\
    \textbf{Software}: Python, C/C++, Java\\
    \textbf{Web/API}: Flask, Django, HTML/CSS, Javascript\\
    \textbf{DevOps}: Git, Docker\\
    \textbf{Low-Level}: x86 assembly, SystemVerilog\\
    \textbf{Robotics}: ROS, OpenCV\\
    \textbf{Research/Math}: \LaTeX\\
    
    
    % -------------------- TEACHING --------------------
    
    \section{Teaching and Mentorship}
    \textbf{Algorithms \& Models of Computation Classroom Assistant}\\ August 2025 -- Present\vspace{-5pt}
    \begin{itemize}[itemsep=-5pt, parsep=5pt, leftmargin=0.25in]
        \item Assist weekly in quiz and exam grading and hold 2-hour-long office hours weekly.
        \item Writing lecture notes on the connection between the WL graph isomorphism test and the expressivity of Graph Neural Networks.
    \end{itemize}
    \textbf{Analog Signal Processing (ECE 210) Tutor (IEEE-HKN)}\\ August 2023 -- May 2024\vspace{-5pt}
    \begin{itemize}[itemsep=-5pt, parsep=5pt, leftmargin=0.25in]
        \item Provided 1-on-1 tutoring, created review slideshows/worksheets, and gave review lectures before midterms.
    \end{itemize}
    
    % -------------------- AWARDS --------------------
    
    \section{Grants \& Awards}
    \begin{itemize}[itemsep=-3pt, parsep=5pt, leftmargin=0.25in]
        \item Awarded grant for undergraduate research from the Office of Undergraduate Research (OUR) \hfill 2025
        \item Awarded James Scholar at University of Illinois Urbana-Champaign \hfill 2023 - 2025
        \item Inducted into Eta Kappa Nu (IEEE-HKN), an ECE Honors Society \hfill 2023
        \item Inducted into Tau Beta Pi, the Engineering Honor Society \hfill 2023
        \item Selected for American Invitational Mathematics Examination (AIME) \hfill 2022
        \item Received the Illinois State Seal of Biliteracy in Spanish \hfill 2020
    \end{itemize}
    
    
    \switchcolumn
    
    % COLUMN 2
    % -------------------- Publications --------------------
    
    \section{Publications \& Conferences}
    \begin{enumerate}[leftmargin=0.25in, itemsep=-2pt]
        \item Hartman, M., \textbf{Jayaraman V.A.}, Choaria, M., Bhimaraju, A., \& Varshney, L.R. (2025). Skip-It? Theoretical Conditions for Layer Skipping in Vision-Language Models. Preprint \href{https://www.arxiv.org/abs/2509.25584}{\underline{arXiv:2509.25584}} \textit{(Under review for ICLR 2026)}
        \item Accepted to the first \href{https://events.ycombinator.com/ai-sus}{\underline{AI Startup School Conference}} (June 2025) hosted by YCombinator for accomplished undergraduate and graduate researchers.
        \item Hartman, M., \textbf{Jayaraman, V.A.}, , Choraria, M., Bhimaraju, A., \& Varshney, L. R. (2025). Unmasking the functionality of early layers in VLMs \textit{(To be presented at eXCV Workshop @ ICCV 2025)}
        \item Wu, X., Hartman, M., \textbf{Jayaraman, V.A.}, \& Varshney, L. R. (2024). SwitchCIT: Switching for Continual Instruction Tuning of Large Language Models. Preprint \href{https://arxiv.org/abs/2407.11780}{\underline{arXiv:2407.11780}}. \textit{(Under review for JSTSP 2025)}
        \item Bernstein, H.C., Bindel, S.R., McKibben, M.A., \& \textbf{Jayaraman, V.A.} (2024). Planning Model based on Projection Methodology Bayesian Discrete Extended (PM2-BDE) \textit{(Under internal review)}
    \end{enumerate}
    
    % -------------------- PROJECTS --------------------
    
    \section{Projects}
    \textbf{RAG Chat-bot for ECE 391 (Computer Systems Engineering)}\\
    \textit{Python, LangChain, dash} \hfill July 2024 -- September 2024\vspace{-5pt}
    \begin{itemize}[leftmargin=0.25in, itemsep=-5pt, parsep=5pt]
        \item Created a RAG chat-bot to search through ECE 391 project documents, citing the document title and page number.
        \item Implemented cutting-edge RAG techniques such as parent-child text splitting and cross-encoder reranking.
    \end{itemize}
    
    \textbf{Operating System}\\
    \textit{C, x86 assembly} \hfill March 2024 -- May 2024\vspace{-5pt}
    \begin{itemize}[leftmargin=0.25in, itemsep=-5pt, parsep=5pt]
        \item Developed kernel for a Linux-based OS with interrupt-based device I/O support.
        \item Implemented radix-2 paging, Round-Robin scheduling, and a file system capable of 4MB files.
        \item Created a Linux shell for user command input of basic Linux shell commands.
    \end{itemize}
    
    \textbf{Named Entity Highlighter | Chrome Extension}\\
    \textit{Javascript, Python, PyTorch, Django} \hfill July 2023 -- August 2023\vspace{-5pt}
    \begin{itemize}[leftmargin=0.25in, itemsep=-5pt, parsep=5pt]
        \item Trained and implemented a state-of-the-art \textbf{Named Entity Recognition (NER) AI model} from scratch.
        \item Developed a Chrome Extension to highlight and provide Wikipedia links to named entities on a web page.
    \end{itemize}
    
\end{paracol}

\end{document}
